\section{Introduction}

Hit-to-lead optimisation turns an initial binder into a lead through rounds of closely related
analogues. The enumeration is not the hard part: a matched-pair tool will hand a project team
thousands of analogues of any hit~\citep{wawer2014whatnext}. The hard part is choosing the five to
make this month and defending the choice, and chemists make that choice by reading the complex ---
a donor left unsatisfied, a subpocket left empty, a carbon pressed against a polar rim. Each
modification is a hypothesis about the pocket, and the measured outcome of the previous round is
what informs the next.

Two literatures address parts of this and meet nowhere in the middle. Pocket-conditioned generative
models place atoms inside a measured binding site and are, by construction, interaction-aware;
some predict contacts or pharmacophores explicitly~\citep{peng2022pocket2mol,guan2023targetdiff,
schneuing2024diffsbdd,feng2024lingo3dmol,jiang2024pocketflow}. None of them says \emph{why}. On the
other side, chemistry language models now propose and justify modifications in natural language,
increasingly with explicit reasoning traces trained from outcome
data~\citep{zhao2024chemdfm,narayanan2025ether0,cavanagh2024smileyllama,liu2026desro,liu2026drugr,
wen2026logos}, and recent work does put pockets and even required residue contacts into an LLM
prompt~\citep{macdougall2026spatial}. What is still missing is the thing a project team actually
consumes: a stated design hypothesis whose claims can be \emph{checked}, one by one, against the
structure the model was shown.

That check is the contribution of this paper. We condition on a per-atom contact table --- for this
ligand, in this pocket, which residue atom sits how far from which ligand atom, with interaction
typing --- and we treat the design intent as a first-class output that is evaluated separately from
the molecule. Because the intent is a sentence that names residues, every clause it asserts can be
matched against the rows of the table the model received, and a sentence that invents a contact can
be caught automatically and counted.

Making the intent evaluable also makes it \emph{replaceable}. We factorise
$p(y,z\mid x,c)=p_R(z\mid x,c)\,p_M(y\mid x,c,z)$, with $x$ the hit, $c$ its contact table, $z$ the
intent and $y$ the analogue, and train $R$ and $M$ as separate adapters with separate losses. The
factorisation is the chain rule, so it costs nothing in expressiveness; what it buys is separate
supervision and separately trained weights for the two factors. It does not buy the bottleneck
itself: a jointly trained model that emits the sentence before the SMILES is already routing the
molecule through the reason in token order, which is why an intent can be teacher-forced into it and
why the comparison between joint and factorised is an empirical question rather than an
architectural one. We
can then intervene on $z$ alone --- substituting the reference intent, a sentence written for
a different pair, or no sentence at all --- and read off how much of the molecule the sentence
decides. Prior reason-then-molecule systems generate the rationale and the structure inside one
model and one trace~\citep{liu2026desron,liu2026drugr,wen2026logos}, so the reason cannot be swapped
and this measurement is unavailable; earlier work factorised \emph{substructure}-to-molecule rather
than sentence-to-molecule~\citep{jin2020rationale}.

Our supervision is also different in kind. Reasoning traces in recent molecular-optimisation work
are manufactured from the outcome: derived by analysing which fragment groups improved a measured
or computed property~\citep{liu2026desro,liu2026drugr}. We instead mine the sentence the authors
themselves wrote --- the hypothesis as stated in the paper, before the assay result was known to the
reader --- from \m{1{,}622} medicinal-chemistry papers, paired with the two compounds the paper
labels and with their potency difference measured \emph{in the same assay} of that same paper. To
our knowledge no prior work extracts the stated design rationale of a medicinal-chemistry
modification from paper text; the existing literature-mining work extracts reactions, procedures and
entities~\citep{chemminer2024,materials2025ie}. We also report, in the same breath, that only
\m{39.3\%} of what we extract states an instruction rather than restating an outcome, which puts the
majority of our own corpus in the same category as the methods we are distinguishing ourselves from.

\paragraph{What we find, including what we did not want to find.}
\begin{enumerate}\itemsep2pt
\item \textbf{An intent sentence behaves like an executable specification.} Given the reference
intent, $M$ reproduces the analogue made at \m{0.919} mean Tanimoto, median \m{1.000}, on a
molecule-level holdout; against a sentence belonging to another pair, \m{0.541} ($+0.379$ paired,
$p=5\times10^{-50}$). A \emph{wrong} sentence is worse than no sentence (\m{0.147} below,
$p=5\times10^{-40}$) and costs validity (\m{100\%}$\rightarrow$\m{91.5\%}). The model follows the
instruction; it does not merely decorate.
\item \textbf{The no-intent arm is not a baseline.} With no sentence, $M$ returns the input
unchanged in \m{100\%} of rows. Every property delta is exactly zero and its Tanimoto to the
reference is the \emph{input's} own similarity. We report it as a diagnostic --- zero intent, zero
edit --- and use the mismatched arm as the honest control.
\item \textbf{The pocket governs the prose, not the molecule.} Withholding the contact table takes
groundedness from \m{1.000} to \m{0.432} and residue specificity from \m{0.388} to \m{0.121}, and
leaves similarity to the reference untouched (\m{0.679} vs.\ \m{0.680}). A deliberately wrong pocket
costs \m{0.013} Tanimoto at cohort scale. We therefore claim \emph{grounded attribution}, not
pocket-aware generation, and we identify the likely mechanism: the fields that would describe
\emph{where a group could go} --- open room, slot shape, exit vectors --- are computable but are
excluded from the model's input because, computed only at the edit site, they leaked it.
\item \textbf{The reason model is the bottleneck.} $R\!\rightarrow\!M$ reaches \m{0.570}, retaining
\m{8.7\%} of the oracle-over-mismatch gap, and $R$'s intents are repetitive (uniqueness \m{0.238}).
The factorisation is validated as an upper bound and as a product claim --- a chemist can type the
intent today --- not as an end-to-end gain.
\item \textbf{Property-directed intents move the property, unevenly.} Polarity ($+3.62$ TPSA,
$n=259$), lipophilicity ($-0.63$ cLogP, $n=146$) and size-down ($-31.2$\,Da, $n=75$) move as asked
and are flat or reversed under a mismatched intent; size-up is flat and rigidity moves the wrong
way.
\end{enumerate}

We also publish the corpus defects, the metrics we refuse to report, and two numbers from this
programme that were withdrawn after failing to reproduce. The cheapest credibility in a paper of
this kind is naming your own retractions.
